% !TEX program = xelatex
% Ficha de comandos: uma por grupo. O grupo escreve os comandos ANTES de ir ao switch;
% o monitor confere e assina.
\documentclass[11pt]{article}
\usepackage[brazil]{babel}
\usepackage[sfdefault,lf]{FiraSans}
\usepackage{FiraMono}
\usepackage[a4paper,margin=1.6cm]{geometry}
\usepackage[table]{xcolor}
\usepackage{array,tabularx,booktabs}
\definecolor{Navy}{HTML}{111827}
\definecolor{Ciano}{HTML}{0B7F7D}
\definecolor{Cinza}{HTML}{5B6677}
\definecolor{RuleGray}{HTML}{C9D1DB}
\pagestyle{empty}
\setlength{\parindent}{0pt}
\arrayrulecolor{RuleGray}
\newcommand{\cab}[1]{\textcolor{white}{\bfseries #1}}
\newcommand{\linhas}[1]{\par\foreach\i in {1,...,#1}{\rule{\linewidth}{0.4pt}\par\vspace{0.46cm}}}
\usepackage{pgffor}
\newcommand{\tarefa}[2]{\vspace{0.35cm}{\color{Ciano}\bfseries\large Tarefa #1:} {\bfseries #2}\par\vspace{0.15cm}}
\newcommand{\visto}{\hfill{\small\color{Cinza}Visto do monitor: \rule{3.2cm}{0.4pt}}\par}

\begin{document}

{\color{Navy}\bfseries\LARGE Ficha de comandos}\hfill{\large Grupo: \rule{1.5cm}{0.4pt}\quad Portas: \rule{2cm}{0.4pt}}

\vspace{0.1cm}
{\small\color{Cinza}Escrevam os comandos aqui \textbf{antes} do seu turno no switch. Sem o visto do monitor, não digite.}

\tarefa{1}{mapear a célula}
\renewcommand{\arraystretch}{1.45}
\begin{tabularx}{\textwidth}{|c|X|X|X|}
\rowcolor{Navy}\cab{Porta} & \cab{Equipamento} & \cab{MAC} & \cab{IP} \\ \hline
\hspace{1.2cm} & & & \\ \hline
 & & & \\ \hline
 & & & \\ \hline
 & & & \\ \hline
 & & & \\ \hline
\end{tabularx}

\tarefa{2}{isolar a célula (VLAN)}
\linhas{5}
\visto

\tarefa{3}{monitorar (mirror para a porta 35)}
\linhas{2}
{\small Para desfazer no fim do turno:}
\linhas{1}
\visto

\tarefa{4}{endurecer}
\linhas{4}
\visto

\vspace{0.2cm}
{\color{Ciano}\bfseries Provas que o grupo apresenta}
\begin{itemize}\setlength{\itemsep}{0pt}\small
\item PUT/GET funcionando entre os dois CLPs da célula, antes e depois da VLAN
\item \texttt{ping} para o CLP da própria célula responde; para o de outra célula, não
\item Print do Wireshark com um pacote S7comm explicado: quem lê, quem escreve, o quê, e onde estaria a senha
\item Notebook plugado no lugar do CLP sem comunicação; porta livre sem link
\end{itemize}

\end{document}
